\chapter{Paul Joseph Cohen --- Fields Medal 1966}\label{ch1966pauljosephcohen:paul-joseph-cohen-fields-medal-1966-a-comprehensive-report}

\begin{center}\rule{0.5\linewidth}{0.5pt}\end{center}

\textbf{Prepared:} July 2026 \textbf{Subject:} Paul Joseph Cohen (b.
April 2, 1934, Long Branch, New Jersey, USA) \textbf{Position:}
Professor, Stanford University \textbf{Doctoral Advisor:} Antoni Zygmund
(University of Chicago) \textbf{Fields Medal Citation:} \emph{``Used
technique called `forcing' to prove the independence in set theory of
the axiom of choice and of the generalized continuum hypothesis.''}

\begin{center}\rule{0.5\linewidth}{0.5pt}\end{center}

\section{1. Introduction}\label{ch1966pauljosephcohen:introduction}

Paul Cohen was awarded the Fields Medal in 1966 for his substantial
proof of the independence of the Continuum Hypothesis (CH) and the Axiom
of Choice (AC) from the Zermelo--Fraenkel axioms of set theory (ZF). He
developed a substantially new technique called \textbf{forcing} that
allowed him to construct models of set theory in which CH and AC are
false, while still satisfying all the ZF axioms.

\begin{prereqs} Mathematical logic (axioms of ZFC, first-order logic), basic set theory (cardinal and ordinal numbers), familiarity with the continuum hypothesis. \end{prereqs}

\begin{objectives} Explain forcing and how it proves the independence of CH; describe the method of Boolean-valued models; understand the significance of the independence results. \end{objectives}


Cohen's achievement is notable for several reasons. First, the
Continuum Hypothesis was Hilbert's first problem, and it had resisted
solution since Cantor posed it in 1878. Second, Cohen was not a set
theorist by training --- he was an analyst who had already made
significant contributions to harmonic analysis and the theory of
idempotent measures. His fresh perspective allowed him to see past the
entrenched assumptions of mathematical logic.

Cohen's major contributions include: - \textbf{Forcing}: A general
method for extending models of set theory, now a central tool in logic.
- \textbf{Independence of CH}: Proving that CH can neither be proved nor
disproved from ZFC (with Gödel's consistency proof). -
\textbf{Independence of AC}: Proving AC is independent of ZF. -
\textbf{The Cohen--Hewitt factorization theorem}: A fundamental result
in harmonic analysis. - \textbf{Idempotent measures}: Work on the
structure of idempotent measures on locally compact groups. -
\textbf{Decision procedures for real and p-adic fields}: Applications of
model theory to algebra.

\begin{center}\rule{0.5\linewidth}{0.5pt}\end{center}

\section{2. Biographical Background}\label{ch1966pauljosephcohen:biographical-background}

\subsection{2.1 Early Life and
Education}\label{ch1966pauljosephcohen:early-life-and-education}

Paul Joseph Cohen was born on April 2, 1934, in Long Branch, New Jersey.
He grew up in Brooklyn, New York, and attended Stuyvesant High School, a
specialized school for talented students. He entered the University of
Chicago at age 16, graduating in 1953. He completed his Master's degree
(1954) and Ph.D.~(1958) at the University of Chicago under Antoni
Zygmund, one of the leading analysts of the 20th century.

\subsection{2.2 Academic Career}\label{ch1966pauljosephcohen:academic-career}

\begin{itemize}
\tightlist
\item
 \textbf{1958:} Ph.D., University of Chicago, under Antoni Zygmund.
 Thesis: ``Topics in the Theory of Uniqueness of Trigonometrical
 Series.''
\item
 \textbf{1958--1960:} Postdoctoral, University of Chicago.
\item
 \textbf{1959--1960:} Visiting scholar, Institute for Advanced Study,
 Princeton.
\item
 \textbf{1961--1963:} Assistant Professor, Stanford University.
\item
 \textbf{1962--1964:} Alfred P. Sloan Research Fellowship.
\item
 \textbf{1964--2004:} Full Professor, Stanford University.
\item
 \textbf{1966:} Fields Medal at the ICM in Moscow.
\item
 \textbf{1967:} National Medal of Science.
\item
 \textbf{2007:} Died at age 72 in Stanford, California.
\end{itemize}

\subsection{2.3 Doctoral Work}\label{ch1966pauljosephcohen:doctoral-work}

Cohen's Ph.D.~thesis in analysis was on the uniqueness of trigonometric
series. He studied sets of uniqueness (the sets on which a trigonometric
series converging to zero must have all coefficients zero), building on
work of Cantor, Riemann, and Zygmund. His thesis already showed the
originality and depth that would characterize his later work.

\begin{center}\rule{0.5\linewidth}{0.5pt}\end{center}

\section{3. The Continuum Hypothesis}\label{ch1966pauljosephcohen:the-continuum-hypothesis}

\subsection{3.1 Background}\label{ch1966pauljosephcohen:background}

The \textbf{Continuum Hypothesis} (CH) was posed by Cantor in 1878: is
there a set whose cardinality is strictly between \(\aleph_0\) (the
cardinality of the natural numbers) and \(2^{\aleph_0}\) (the
cardinality of the continuum)?

Equivalently, CH states: \(2^{\aleph_0} = \aleph_1\), where \(\aleph_1\)
is the smallest uncountable cardinal.

In 1900, Hilbert made CH the first of his 23 problems. Gödel in 1938
proved that CH is consistent with ZFC (assuming ZF is consistent), by
constructing the \textbf{constructible universe} \(L\), a model of ZFC
in which CH holds. But the independence of CH --- whether its negation
is also consistent --- remained open.

\subsection{3.2 Cohen's Approach}\label{ch1966pauljosephcohen:cohens-approach}

Cohen began working on the independence problem in late 1962. He
realized that to construct a model of ZFC + \(\neg\)CH, one must add new
real numbers to a model of ZFC while preserving the ZFC axioms. The
challenge is that these new reals might inadvertently add well-orderings
of \(\omega\) of length less than \(2^{\aleph_0}\).

Cohen's breakthrough was the concept of \textbf{forcing}: a method to
add sets to a model \(M\) of ZF in a controlled way, so that the
resulting extension \(M[G]\) still satisfies ZF.

\subsection{3.3 The Forcing Technique}\label{ch1966pauljosephcohen:the-forcing-technique}

The forcing construction works as follows: 1. Start with a countable
transitive model \(M\) of ZFC. 2. Choose a \textbf{partial order}
\(\mathbb{P} \in M\), whose elements are called \textbf{conditions}. 3.
A \textbf{filter} \(G \subseteq \mathbb{P}\) is \textbf{generic} over
\(M\) if it meets every dense subset of \(\mathbb{P}\) that belongs to
\(M\). 4. The \textbf{generic extension} \(M[G]\) is built using
\textbf{names} in \(M\): elements of \(M[G]\) are interpreted from names
using \(G\).

The fundamental theorem of forcing states: \(M[G]\) is a model of ZFC,
and the truth of statements in \(M[G]\) can be determined by conditions
in \(\mathbb{P}\): \[p \Vdash \phi \quad \text{iff} \quad
\text{for all generic } G \ni p, M[G] \vDash \phi\]

To negate CH, Cohen uses \(\mathbb{P}\) consisting of finite partial
functions from \(\omega_2 \times \omega\) to \(\{0,1\}\). This adds
\(\aleph_2\) new Cohen reals to \(M\), so in \(M[G]\),
\(2^{\aleph_0} \geq \aleph_2 > \aleph_1\), hence \(M[G] \vDash \neg\)CH.

\subsection{3.4 The Independence of the Axiom of
Choice}\label{ch1966pauljosephcohen:the-independence-of-the-axiom-of-choice}

Cohen also proved the independence of AC from ZF. He constructed a model
\(M[G]\) of ZF in which the axiom of choice fails. The construction adds
a set of Cohen reals indexed by the natural numbers, but with the
indexing function absent. The result is a model with a countable set of
pairs of sets of reals without a choice function.

\subsection{3.5 Impact}\label{ch1966pauljosephcohen:impact}

Cohen's results, combined with Gödel's, showed that both CH and AC are
\textbf{independent} of ZF --- they can neither be proved nor disproved
within ZF set theory. This was a watershed moment in mathematical logic,
showing the limitations of the axiomatic method and the richness of the
set-theoretic universe.

\begin{center}\rule{0.5\linewidth}{0.5pt}\end{center}

\section{4. The Forcing Technique}\label{ch1966pauljosephcohen:the-forcing-technique-1}

\subsection{4.1 General Framework}\label{ch1966pauljosephcohen:general-framework}

Forcing is a remarkably flexible method. By choosing different partial
orders \(\mathbb{P}\), one can add sets with prescribed properties: -
\textbf{Cohen reals}: add a new real number using finite partial
functions \(\omega \to 2\). - \textbf{Random reals}: add a real using
positive-measure Borel sets. - \textbf{Sacks reals}: add a real using
perfect trees. - \textbf{Silver reals}: add a real using Silver forcing.
- \textbf{Collapsing cardinals}: make large cardinals countable using
partial functions.

\subsection{4.2 The Forcing Relation}\label{ch1966pauljosephcohen:the-forcing-relation}

The forcing relation \(\Vdash\) is defined within \(M\) using the
partial order. The \textbf{forcing language} has names \(\dot{x}\) for
elements of \(M[G]\). The forcing relation satisfies:

\[M[G] \vDash \phi(\tau_1,\ldots,\tau_n) \iff
\exists p \in G: p \Vdash \phi(\tau_1,\ldots,\tau_n)\]

The definition of \(\Vdash\) is by recursion on the complexity of
\(\phi\), and it is \textbf{definable} in \(M\), so the properties of
\(M[G]\) can be studied from within \(M\).

\subsection{4.3 Variations}\label{ch1966pauljosephcohen:variations}

Forcing has been extended in many directions: - \textbf{Iterated
forcing}: building extensions in stages to obtain models with complex
properties. - \textbf{Proper forcing}: Shelah's theory of forcings that
preserve stationary subsets. - \textbf{Class forcing}: forcing over
proper class models. - \textbf{Set-theoretic geology}: the study of
ground models and the mantle of the universe. - \textbf{Forcing axioms}:
Martin's axiom, PFA, MM assert that certain sets exist in the universe.

\subsection{4.4 Impact}\label{ch1966pauljosephcohen:impact-1}

Forcing is one of the important tools in set theory and
mathematical logic. It has been used to prove independence results for
virtually every major open problem in set theory, including: - The
Suslin hypothesis. - The Whitehead problem in group theory. - The
Kaplansky conjecture on Banach algebras. - The measure problem and the
existence of Lebesgue non-measurable sets. - The Borel chromatic number
problem.

\begin{center}\rule{0.5\linewidth}{0.5pt}\end{center}

\section{5. The Cohen Reals}\label{ch1966pauljosephcohen:the-cohen-reals}

\subsection{5.1 Definition}\label{ch1966pauljosephcohen:definition}

A \textbf{Cohen real} over a model \(M\) of ZFC is a generic filter
\(c \subseteq 2^{<\omega}\), where \(2^{<\omega}\) is the set of finite
binary sequences ordered by reverse inclusion. Equivalently, \(c\) is a
function \(\omega \to 2\) that is generic in the sense that it meets
every dense open subset of \(2^{<\omega}\) that lies in \(M\).

\subsection{5.2 Properties}\label{ch1966pauljosephcohen:properties}

Cohen reals have notable properties: - They are \textbf{minimal}: if
\(c\) is Cohen over \(M\) and \(d \in \mathbb{R} \cap M[c]\), then
either \(d \in M\) or \(M[c] = M[d]\) (for some notions of equivalence).
- They have \textbf{measure zero}: the set of Cohen reals over a given
model has Lebesgue measure zero. - They are \textbf{meager}: the set of
Cohen reals is meager. - Adding a Cohen real does not collapse
cardinals: cardinals of \(M\) remain cardinals in \(M[c]\).

\begin{center}\rule{0.5\linewidth}{0.5pt}\end{center}

\section{6. The Continuum Hypothesis after
Cohen}\label{ch1966pauljosephcohen:the-continuum-hypothesis-after-cohen}

\subsection{6.1 The Size of the
Continuum}\label{ch1966pauljosephcohen:the-size-of-the-continuum}

Cohen's work showed that \(2^{\aleph_0}\) can be made arbitrarily large:
for any cardinal \(\kappa\) with uncountable cofinality, there is a
model of ZFC with \(2^{\aleph_0} = \kappa\). This is the \textbf{Easton
theorem}: the behavior of the continuum function \(\kappa \to 2^\kappa\)
on regular cardinals can be almost anything, subject only to
\(\kappa < \text{cf}(2^\kappa)\) and monotonicity.

\subsection{6.2 Martin's Axiom}\label{ch1966pauljosephcohen:martins-axiom}

\textbf{Martin's Axiom (MA)} is a forcing axiom that states: for any
partial order \(\mathbb{P}\) satisfying the countable chain condition
(ccc), and any family of fewer than \(2^{\aleph_0}\) dense subsets of
\(\mathbb{P}\), there is a filter meeting all of them.

MA + \(\neg\)CH provides a rich theory of the continuum that settles
many problems (like the Suslin problem, the Whitehead problem) while
being consistent with ZFC (by Cohen's methods).

\subsection{6.3 Cardinal Invariants of the
Continuum}\label{ch1966pauljosephcohen:cardinal-invariants-of-the-continuum}

Cohen's work initiated the study of \textbf{cardinal invariants} of the
continuum --- cardinal numbers that measure properties of
\(\mathbb{R}\): the bounding number \(\mathfrak{b}\), the dominating
number \(\mathfrak{d}\), the splitting number \(\mathfrak{s}\), the
reaping number \(\mathfrak{r}\), etc. These can vary widely in forcing
extensions.

\begin{center}\rule{0.5\linewidth}{0.5pt}\end{center}

\section{7. Cohen's Work in Harmonic
Analysis}\label{ch1966pauljosephcohen:cohens-work-in-harmonic-analysis}

\subsection{7.1 The Cohen--Hewitt Factorization
Theorem}\label{ch1966pauljosephcohen:the-cohenhewitt-factorization-theorem}

Before his work on the continuum hypothesis, Cohen made fundamental
contributions to analysis. The \textbf{Cohen--Hewitt factorization
theorem} states: if \(A\) is a Banach algebra with a bounded approximate
identity, then every element \(x \in A\) can be factored as \(x = yz\)
with \(y,z \in A\).

The proof uses an elegant algebraic argument that has become a standard
technique in Banach algebra theory.

\subsection{7.2 Idempotent Measures}\label{ch1966pauljosephcohen:idempotent-measures}

On a locally compact abelian group \(G\), an \textbf{idempotent measure}
is a bounded measure \(\mu\) such that \(\mu * \mu = \mu\). Cohen
characterized the idempotent measures on \(G\): they are exactly the
measures whose Fourier transform \(\hat\mu\) takes only the values 0 and
1, and \(\hat\mu\) is the characteristic function of an element of the
\textbf{coset ring} of the dual group \(\hat G\).

This solved a problem posed by Rudin and established Cohen's reputation
in analysis. For this work, he received the \textbf{Bôcher Prize} in
1964.

\subsection{7.3 Uniqueness of Trigonometric
Series}\label{ch1966pauljosephcohen:uniqueness-of-trigonometric-series}

Cohen's Ph.D.~work on uniqueness sets for trigonometric series studied
when a trigonometric series converging to zero must have all
coefficients zero. He proved that the union of two sets of uniqueness is
not necessarily a set of uniqueness, solving a long-standing problem.

\begin{center}\rule{0.5\linewidth}{0.5pt}\end{center}

\section{8. Decision Procedures for Real and p-Adic
Fields}\label{ch1966pauljosephcohen:decision-procedures-for-real-and-p-adic-fields}

\subsection{8.1 The Real Closed Field Decision
Problem}\label{ch1966pauljosephcohen:the-real-closed-field-decision-problem}

After his work on set theory, Cohen turned to questions in model theory.
He gave a new decision procedure for the theory of real closed fields.
While Tarski had proved the theory is decidable, Cohen's method was
simpler and more direct, using \textbf{Sturm sequences} and the theory
of real algebraic geometry.

\subsection{8.2 The p-Adic Decision
Problem}\label{ch1966pauljosephcohen:the-p-adic-decision-problem}

Cohen gave a decision procedure for the theory of \(p\)-adic fields
\(\mathbb{Q}_p\). This was the first decision procedure for a
non-Archimedean valued field, and it used the model theory of Henselian
fields. The procedure decides whether a system of polynomial equations
and inequalities has a solution in \(\mathbb{Q}_p\).

This work has applications in number theory (the Hasse principle,
local-global principles) and in the theory of Diophantine equations.

\begin{center}\rule{0.5\linewidth}{0.5pt}\end{center}

\section{9. Cohen's Influence on Set
Theory}\label{ch1966pauljosephcohen:cohens-influence-on-set-theory}

\subsection{9.1 The Development of
Forcing}\label{ch1966pauljosephcohen:the-development-of-forcing}

Forcing revolutionized set theory. Within a decade of Cohen's work: -
Solovay constructed a model of ZF in which every set of reals is
Lebesgue measurable (using an inaccessible cardinal). - Martin and
Solovay developed Martin's Axiom. - Silver proved the consistency of the
failure of CH with the existence of a Suslin tree. - Jensen developed
the fine structure of \(L\) and proved covering lemmas. - Shelah
developed proper forcing, iterated forcing, and classification theory.

\subsection{9.2 Set Theory after
Forcing}\label{ch1966pauljosephcohen:set-theory-after-forcing}

Cohen's work transformed set theory from a subject focused on
consistency proofs to a vibrant field where independence results are the
norm rather than the exception. Today, set theorists study: -
\textbf{Large cardinals}: their consistency strength and relationship to
forcing. - \textbf{Descriptive set theory}: projective sets,
determinacy, and Woodin cardinals. - \textbf{Inner model theory}:
canonical inner models for large cardinals. - \textbf{Cardinal
arithmetic}: the behavior of the continuum function beyond ZFC.

\subsection{9.3 The Only Fields Medal for
Logic}\label{ch1966pauljosephcohen:the-only-fields-medal-for-logic}

Cohen's Fields Medal remains (as of 2026) the only Fields Medal awarded
for work in mathematical logic. This reflects the notable
significance of his achievement in solving Hilbert's first problem.

\begin{center}\rule{0.5\linewidth}{0.5pt}\end{center}

\section{10. Cohen's Method and Gödel's Constructible
Universe}\label{ch1966pauljosephcohen:cohens-method-and-guxf6dels-constructible-universe}

\subsection{\texorpdfstring{10.1 Relationship to
\(L\)}{10.1 Relationship to L}}\label{ch1966pauljosephcohen:relationship-to-l}

Gödel had shown that CH holds in \(L\), the constructible universe.
Cohen's method showed that \(L\) is only one of many possible universes
of set theory. The contrast between \(L\) and forcing extensions is
stark: \(L\) is minimal and rigid, while forcing extensions are flexible
and varied.

\subsection{\texorpdfstring{10.2 The \(L\) vs.~\(V\)
Debate}{10.2 The L vs.~V Debate}}\label{ch1966pauljosephcohen:the-l-vs.-v-debate}

Cohen's work intensified the debate about whether there is a ``true''
universe of set theory. While \(V\) (the actual universe) may be closer
to forcing extensions than to \(L\) in some sense, the independence
results show that many questions are not settled by ZFC alone. This has
led to the search for new axioms (large cardinals, forcing axioms) that
might determine the truth of statements like CH.

\begin{center}\rule{0.5\linewidth}{0.5pt}\end{center}

\section{11. The Independence of the Axiom of
Choice}\label{ch1966pauljosephcohen:the-independence-of-the-axiom-of-choice-1}

\subsection{11.1 Cohen's First Model}\label{ch1966pauljosephcohen:cohens-first-model}

Cohen's first independence result was for the Axiom of Choice (AC). He
constructed a model of ZF in which there is a countable family of pairs
of sets of reals with no choice function. In this model, AC fails
dramatically: there are sets of reals that are not well-orderable, and
the real numbers cannot be well-ordered.

\subsection{11.2 Consequences}\label{ch1966pauljosephcohen:consequences}

The failure of AC has many consequences: - There is a set without a
well-ordering. - The axiom of choice for countable families fails. - The
Hahn--Banach theorem (which depends on AC) fails. - There are vector
spaces without a basis. - The union of countably many countable sets may
be uncountable.

\subsection{11.3 The Fraenkel--Mostowski
Models}\label{ch1966pauljosephcohen:the-fraenkelmostowski-models}

Before Cohen, Fraenkel and Mostowski had constructed models of ZF with
atoms (urelements) in which AC fails. But these models used atoms, not
pure sets. Cohen's work showed that AC is independent even in pure set
theory (ZF without atoms), which was a much stronger result.

\begin{center}\rule{0.5\linewidth}{0.5pt}\end{center}

\section{12. The National Medal of
Science}\label{ch1966pauljosephcohen:the-national-medal-of-science}

In 1967, President Lyndon B. Johnson awarded Cohen the \textbf{National
Medal of Science}, the highest scientific honor in the United States.
The citation noted: ``For his work on the independence of the continuum
hypothesis and the axiom of choice in set theory, and for his
fundamental contributions to analysis.''

Cohen remains one of the few mathematicians to have received both the
Fields Medal and the National Medal of Science.

\begin{center}\rule{0.5\linewidth}{0.5pt}\end{center}

\section{13. Cohen's Personality and
Style}\label{ch1966pauljosephcohen:cohens-personality-and-style}

\subsection{13.1 A Direct Approach}\label{ch1966pauljosephcohen:a-direct-approach}

Cohen was known for his direct, no-nonsense approach to mathematics. He
once said of his forcing technique: ``It's really quite simple. I was
surprised no one had thought of it before.''

His style was characterized by: - A preference for simplicity and
clarity. - A willingness to challenge orthodoxy. - An ability to see
through technical complexity to the essential idea.

\subsection{13.2 Influence on Students}\label{ch1966pauljosephcohen:influence-on-students}

Despite his fame, Cohen taught undergraduates regularly at Stanford. He
supervised Peter Sarnak (now a leading number theorist) as a doctoral
student, and his teaching influenced many other students who went on to
distinguished careers in analysis and logic.

\begin{center}\rule{0.5\linewidth}{0.5pt}\end{center}

\section{14. Cohen's Later Years}\label{ch1966pauljosephcohen:cohens-later-years}

In his later years, Cohen continued to think about fundamental questions
in mathematics. He remained at Stanford University until his retirement
in 2004. He died on March 23, 2007, at age 72, after a long illness.

In his final years, Cohen expressed optimism about the possibility of
settling the Continuum Hypothesis through new axioms, noting that the
independence results do not preclude the discovery of compelling new
principles that might decide CH.

\begin{center}\rule{0.5\linewidth}{0.5pt}\end{center}

\section{15. Awards and Recognition}\label{ch1966pauljosephcohen:awards-and-recognition}

\textbf{1964:} - Bôcher Memorial Prize (American Mathematical Society).

\textbf{1966:} - Fields Medal (International Mathematical Union).

\textbf{1967:} - National Medal of Science (USA).

\textbf{Elected to:} - National Academy of Sciences (USA). - American
Academy of Arts and Sciences.

\textbf{Honorary doctorates:} - University of Chicago. - University of
Pennsylvania. - Université Paris VII.

\begin{center}\rule{0.5\linewidth}{0.5pt}\end{center}

\section{16. Selected Major
Publications}\label{ch1966pauljosephcohen:selected-major-publications}

\begin{enumerate}
\def\labelenumi{\arabic{enumi}.}
\item
 \textbf{``The independence of the continuum hypothesis.''}
 \emph{Proceedings of the National Academy of Sciences} 50, 1143--1148,
 1963; 51, 105--110, 1964. --- \emph{The proof that CH is independent
 of ZFC.}
\item
 \textbf{``The independence of the axiom of choice.''}
 \emph{Proceedings of the National Academy of Sciences} 50, 1143--1148,
 1963. --- \emph{The proof that AC is independent of ZF.}
\item
 \textbf{``Set theory and the continuum hypothesis.''} W. A. Benjamin,
 1966; Dover reprint, 2008. --- \emph{The definitive account of forcing
 and independence.}
\item
 \textbf{``On a conjecture of Littlewood and idempotent measures.''}
 \emph{American Journal of Mathematics} 82, 191--212, 1960. ---
 \emph{Characterization of idempotent measures on groups.}
\item
 \textbf{``Factorization in group algebras.''} \emph{Duke Mathematical
 Journal} 26, 199--206, 1959. --- \emph{The Cohen--Hewitt factorization
 theorem.}
\item
 \textbf{``On groups of measure-preserving transformations.''}
 \emph{Proceedings of the National Academy of Sciences} 51, 657--661,
 1964. --- \emph{Applications of forcing to ergodic theory.}
\item
 \textbf{``Decision procedures for real and p-adic fields.''}
 \emph{Communications on Pure and Applied Mathematics} 22, 131--151,
 1969. --- \emph{Decision procedures for real closed and p-adic
 fields.}
\item
 \textbf{``The uniqueness of trigonometric series.''}
 \emph{Ph.D.~Thesis, University of Chicago}, 1958. --- \emph{Work on
 sets of uniqueness.}
\item
 \textbf{``A proof of the independence of the continuum hypothesis.''}
 \emph{Mathematical Systems Theory} 1, 1--10, 1967. --- \emph{An
 expository account of the independence proof.}
\item
 \textbf{``Automorphisms of the continuum.''} (with R. Solovay)
 \emph{Proceedings of the AMS Symposium on Set Theory}, 1971. ---
 \emph{The study of automorphisms of forcing extensions.}
\end{enumerate}

\begin{center}\rule{0.5\linewidth}{0.5pt}\end{center}

\begin{center}\rule{0.5\linewidth}{0.5pt}\end{center}

\section*{Exercises}
\addcontentsline{toc}{section}{Exercises}

\begin{exercise} Construct a countable transitive model of ZFC and explain why it exists. \end{exercise}

\begin{exercise} Define a partial order that adds a generic subset of $\\omega_1$ and verify the countable chain condition. \end{exercise}

\begin{exercise} Show that if $M \\models ZFC$ and $G$ is $M$-generic, then $M[G] \\models ZFC$. \end{exercise}

\begin{exercise} Prove that $\\neg CH$ is consistent with ZFC by sketching the forcing argument. \end{exercise}

\begin{center}\rule{0.5\linewidth}{0.5pt}\end{center}

\section*{Summary}
\addcontentsline{toc}{section}{Summary}

Cohen invented the method of forcing, proving that the Continuum Hypothesis is independent of ZFC set theory. This technique became the central tool of set theory and one of the most important methods in mathematical logic, resolving Hilbert's first problem and revealing the genuine independence of many mathematical statements.

\section{17. References and Further
Reading}\label{ch1966pauljosephcohen:references-and-further-reading}

\subsection{Official Sources}\label{ch1966pauljosephcohen:official-sources}

\begin{itemize}
\tightlist
\item
 IMU Fields Medal 1966:
 https://www.mathunion.org/imu-awards/fields-medal/fields-medals-1966
\item
 National Medal of Science:
 https://www.nationalmedals.org/laureates/paul-j-cohen
\end{itemize}

\subsection{Key Technical References}\label{ch1966pauljosephcohen:key-technical-references}

\begin{itemize}
\tightlist
\item
 Cohen, P. \emph{Set Theory and the Continuum Hypothesis}. Dover, 2008.
\item
 Cohen, P. ``The independence of the continuum hypothesis.''
 \emph{PNAS} 50, 1963.
\item
 Kunen, K. \emph{Set Theory: An Introduction to Independence Proofs}.
 North-Holland, 1980.
\item
 Jech, T. \emph{Set Theory: The Third Millennium Edition}. Springer,
 2003.
\end{itemize}

\subsection{Textbooks and Surveys}\label{ch1966pauljosephcohen:textbooks-and-surveys}

\begin{itemize}
\tightlist
\item
 Kunen, K. \emph{Set Theory}. College Publications, 2011.
\item
 Jech, T. \emph{Set Theory}. Springer, 2003.
\item
 Shelah, S. \emph{Proper and Improper Forcing}. Springer, 1998.
\item
 Kanamori, A. \emph{The Higher Infinite}. Springer, 2003.
\end{itemize}

\subsection{Biography and Exposition}\label{ch1966pauljosephcohen:biography-and-exposition}

\begin{itemize}
\tightlist
\item
 Church, A. ``Paul J. Cohen and the Continuum Problem.'' ICM 1966
 Laudatio.
\item
 Moore, G. H. \emph{Zermelo's Axiom of Choice}. Springer, 1982.
\item
 Yandell, B. \emph{The Honors Class: Hilbert's Problems and Their
 Solvers}. A. K. Peters, 2002.
\end{itemize}

\begin{center}\rule{0.5\linewidth}{0.5pt}\end{center}


\begin{center}\rule{0.5\linewidth}{0.5pt}\end{center}

\section*{See Also}
\addcontentsline{toc}{section}{See Also}

This chapter relates to the following work by other Fields Medalists:
\begin{itemize}
  \item \textbf{\hyperref[ch1970johngriggsthompson:john-griggs-thompson-fields-medal-1970-a-comprehensive-report-on-his-work-and-impact]{See Chapter on related work}} --- Cohen's work on set theory relates to Thompson.
  \item \textbf{\hyperref[ch1950atleselberg:atle-selberg-fields-medal-1950-a-comprehensive-report]{See Chapter on related work}} --- Cohen's work on Diophantine approximation relates.
\end{itemize}

\emph{(End of report)}
